\documentclass[a4paper,11pt]{article}
\usepackage[utf8]{inputenc}
\usepackage[T1]{fontenc}
\usepackage{lmodern}
\usepackage{amsmath,amssymb,amsthm}
\usepackage{mathtools}
\usepackage{booktabs}
\usepackage{longtable}
\usepackage{array}
\usepackage{hyperref}
\usepackage{graphicx}
\usepackage{float}
\usepackage{geometry}
\geometry{margin=25mm}

\hypersetup{
  colorlinks=true,
  linkcolor=blue,
  urlcolor=blue,
  citecolor=blue
}

\theoremstyle{definition}
\newtheorem{theorem}{Theorem}[section]
\newtheorem{proposition}[theorem]{Proposition}
\newtheorem{lemma}[theorem]{Lemma}
\newtheorem{corollary}[theorem]{Corollary}
\newtheorem{definition}[theorem]{Definition}
\newtheorem{remark}[theorem]{Remark}
\newtheorem{observation}[theorem]{Observation}

\setlength{\parindent}{1em}
\setlength{\parskip}{0.5em}

\providecommand{\tightlist}{\setlength{\itemsep}{0pt}\setlength{\parskip}{0pt}}
\begin{document}

\section{On Dimension Addition to Background Spaces and Subjective Spaces}\label{on-dimension-addition-to-background-spaces-and-subjective-spaces}

\begin{center}\rule{0.5\linewidth}{0.5pt}\end{center}

\textbf{Author}: Noriaki Kihara (WF System Co., Ltd.)

\textbf{Date}: April 2026

\textbf{DOI:} \href{https://doi.org/10.5281/zenodo.19526913}{10.5281/zenodo.19526913}

\begin{center}\rule{0.5\linewidth}{0.5pt}\end{center}

\subsection{Abstract}\label{abstract}

This paper examines the \textbf{addition of new dimensions} to background spaces and subjective spaces with finite curvature, taking as premises the central projection framework established in prior papers {[}1{]}, {[}2{]}, {[}3{]}, and the existence proposition for nested structures in subjective spaces demonstrated in paper {[}4{]}. While paper {[}4{]} established the geometric possibility of nested structures, the conditions under which an added dimension admits a consistent interpretation require separate investigation. Departing from an observation that extends Theorem 3.1 of paper {[}2{]} --- namely, that \textbf{which curvature radius governs an added dimension is interpretation-dependent on the part of the observer} --- this paper shows that the value of the parent space's curvature radius \(R\) (whether finite or \(\infty\)) imposes structural constraints on the types of dimensions that can be added. The scope of this paper is limited to a geometric and logical description of these constraints; no claim is made regarding their physical reality.

\begin{center}\rule{0.5\linewidth}{0.5pt}\end{center}

\subsection{1. Introduction}\label{introduction}

In the central projection framework introduced in prior paper {[}1{]}, a hypersphere \(S^n(R)\) of radius \(R\) and a tangent hyperplane \(\Pi_R\) are defined within an \((n+1)\)-dimensional Euclidean background space \(\mathbb{R}^{n+1}\). The subjective space \(\Pi_R\) is an \(n\)-dimensional space governed by curvature radius \(R\), satisfying the Einstein equation \(G_{\mu\nu} + \Lambda g_{\mu\nu} = 0\), \(\Lambda = 3/R^2\).

Theorem 3.1 of paper {[}2{]} established that the axes of subjective space share a geometric absolute symmetry, and that which axis is interpreted in a particular role (e.g., as the axis corresponding to time) is left to the observer's choice.

Section §4 of paper {[}4{]} showed that the freedom to initialize a new subjective space within any given subjective space is directly derivable from the existing framework (Proposition 4.1: existence of nested structures).

However, Proposition 4.1 of paper {[}4{]} is an existence proposition asserting that such initialization is possible; the conditions under which the initialized structure admits a \textbf{consistent interpretation} require separate examination. The purpose of this paper is to investigate these consistency conditions for interpretations according to the value of the parent space's curvature radius \(R\).

\begin{center}\rule{0.5\linewidth}{0.5pt}\end{center}

\subsection{2. Interpretation-Dependence of Dimension Addition}\label{interpretation-dependence-of-dimension-addition}

\subsubsection{2.1 General Character of Interpretation-Dependence}\label{general-character-of-interpretation-dependence}

In Theorem 3.1 of paper {[}2{]}, the symmetry among axes of subjective space implies that individual interpretations --- such as which axis ``corresponds to time'' --- are left to the observer's choice. This paper extends this interpretation-dependence to the context of dimension addition.

When a new dimension is added to a subjective space or background space, \textbf{which curvature radius governs the new dimension} is interpretation-dependent on the part of the observer. Specifically, both the stance that interprets the new dimension as governed by the existing curvature radius \(R\), and the stance that interprets it as governed by a newly independent curvature radius \(R'\), may in principle be admitted.

Even when the geometric operation of addition is the same, different interpretations yield different structures for what the new dimension signifies. The ``constraints'' described in this paper are not geometric prohibitions but logical constraints arising from interpretive consistency.

\subsubsection{2.2 Two Types of Dimension Addition}\label{two-types-of-dimension-addition}

From the foregoing discussion, the following two types of dimension addition are distinguished.

\textbf{Type I} (addition under existing governance): The case in which the new dimension is interpreted as governed by the existing curvature radius \(R\). In this case, the new dimension is understood as an extension of the existing subjective space and introduces no new curvature structure.

\textbf{Type II} (introduction of a new curvature radius): The case in which the new dimension is interpreted as governed by a newly independent curvature radius \(R'\). In this case, the new dimension corresponds to the initialization of a distinct subjective space \(\Pi_{R'}\).

In what follows, we examine whether each of these two types of addition is admissible depending on the value of the parent space's curvature radius \(R\).

\begin{center}\rule{0.5\linewidth}{0.5pt}\end{center}

\subsection{3. Dimension Addition to a Flat Background Space}\label{dimension-addition-to-a-flat-background-space}

\subsubsection{3.1 Setup}\label{setup}

The background space \(\mathbb{R}^{n+1}\) carries no curvature. It is the starting point of the entire framework and may be regarded as corresponding to the state \(R = \infty\).

We consider the case of adding a new dimension to this background space.

\subsubsection{3.2 Admissible Additions}\label{admissible-additions}

\textbf{Observation 3.1.} The conditions under which a new dimension admits a consistent interpretation with respect to the flat background space \(\mathbb{R}^{n+1}\) are that one of the following is satisfied.

\textbf{(a)} The new dimension is interpreted as governed by \(R_1 = \infty\). That is, the flatness of the background space is maintained after the addition.

\textbf{(b)} The new dimension is interpreted as governed by an existing or new finite curvature radius, and this interpretation is consistent from the perspective of the background space's inhabitants.

Condition (a) corresponds to an operation that increases the dimensionality of the background space, and is equivalent to extending the dimensionality from \((n+1)\) to \((n+2)\) while preserving flatness. In this case, consistency is trivially satisfied both geometrically and interpretively.

Condition (b) is admissible only when the added dimension can be interpreted as ``part of some finite curvature radius'' from the perspective of the background space's inhabitants. When this interpretation does not hold --- that is, when the new dimension is interpreted as introducing a curvature structure independent of the flatness of the background space --- combining with the axis symmetry of Theorem 3.1 of paper {[}2{]} causes the relevant curvature structure to pervade the entire background space, contradicting the flatness of the background space, which is the starting point.

\subsubsection{3.3 Summary}\label{summary}

For a flat background space, Type I addition is always admissible (Condition (a)). Type II addition is admissible only when the structure governed by the new curvature radius can be consistently interpreted from the perspective of the background space's inhabitants (Condition (b)).

\begin{center}\rule{0.5\linewidth}{0.5pt}\end{center}

\subsection{4. Dimension Addition to a Subjective Space with Finite Curvature}\label{dimension-addition-to-a-subjective-space-with-finite-curvature}

\subsubsection{4.1 Setup}\label{setup-1}

The subjective space \(\Pi_R\) is an \(n\)-dimensional space governed by a finite curvature radius \(R < \infty\). We consider the case of adding a new dimension to this subjective space.

\subsubsection{4.2 Type I Addition: Under Governance of the Existing Curvature Radius}\label{type-i-addition-under-governance-of-the-existing-curvature-radius}

\textbf{Observation 4.1.} For a subjective space \(\Pi_R\), when a new dimension is interpreted as governed by the existing curvature radius \(R\), this Type I addition is \textbf{executable without restriction}.

Since the added dimension shares the same \(R\) as the existing subjective space \(\Pi_R\), it introduces no new curvature structure. By the axis symmetry of Theorem 3.1 of paper {[}2{]}, the distinction between an existing axis and an added axis is a matter of the observer's interpretation and is geometrically symmetric. Therefore, there is in principle no constraint on the number of dimensions that can be added.

\subsubsection{4.3 Type II Addition: Introduction of a New Curvature Radius}\label{type-ii-addition-introduction-of-a-new-curvature-radius}

\textbf{Observation 4.2.} For a subjective space \(\Pi_R\), when a new dimension is interpreted as governed by a newly independent curvature radius \(R' \neq R\), its initialization is geometrically possible by Proposition 4.1 of paper {[}4{]}. However, the addition of this new dimension \textbf{alone} results in a degenerate structure for the initialized subjective space \(\Pi_{R'}\).

Specifically, the axis corresponding to the new curvature radius \(R'\) itself corresponds to the external direction of \(\Pi_{R'}\) and is not perceived by the inhabitants of \(\Pi_{R'}\). Accordingly, the dimensions visible to the inhabitants of \(\Pi_{R'}\) must be constituted by additional dimensions interpreted as governed by \(R'\).

If these additional dimensions are not added simultaneously, the number of dimensions of \(\Pi_{R'}\) as seen by the inhabitants of \(\Pi_{R'}\) is \(0\). That is, \(\Pi_{R'}\) becomes a \textbf{zero-dimensional subjective space}.

\textbf{Observation 4.3.} Therefore, for a Type II addition to constitute a non-degenerate subjective space \(\Pi_{R'}\), it is necessary to \textbf{simultaneously add} the axis corresponding to the new curvature radius \(R'\) and the additional dimensions interpreted as governed by \(R'\). These additional dimensions may be added in any number under \(R'\)'s governance, in the same form as Type I additions.

\subsubsection{4.4 Summary}\label{summary-1}

For a subjective space with finite curvature, Type I additions may be executed in any number. Type II additions result in a zero-dimensional new subjective space unless additional dimensions governed by the new curvature radius \(R'\) are added simultaneously.

\begin{center}\rule{0.5\linewidth}{0.5pt}\end{center}

\subsection{5. Discussion}\label{discussion}

In this paper, we have described the interpretive consistency conditions for dimension addition based on the framework of prior papers {[}1{]}, {[}2{]}, {[}3{]}, {[}4{]}.

The claims of this paper are limited to the following.

\begin{enumerate}
\def\labelenumi{\arabic{enumi}.}
\tightlist
\item
  Types of dimension addition are distinguished, in an interpretation-dependent manner, into two kinds: Type I (addition under existing governance) and Type II (introduction of a new curvature radius).
\item
  For a flat background space, Type I is always admissible, and Type II is admissible only when the interpretation is consistent from the perspective of the background space's inhabitants.
\item
  For a subjective space with finite curvature, Type I is admissible without restriction, and Type II results in a zero-dimensional new subjective space unless additional dimensions governed by the new curvature radius are added simultaneously.
\end{enumerate}

These constraints are not absolute geometric prohibitions but logical constraints arising from the consistency of the observer's interpretation. In this sense, the constraints of this paper lie on the natural extension of the axis symmetry of Theorem 3.1 in paper {[}2{]} --- that is, the fact that the choice of which axis corresponds to the time axis depends on the observer's interpretation.

This paper makes no claim regarding the physical reality of these constraints. Nor does it venture into the physical significance of a new subjective space that becomes zero-dimensional, or the physical implications of the condition ``simultaneously add the governed dimensions'' in Type II addition. These questions lie outside geometry and are left to the reader's judgment.

\begin{center}\rule{0.5\linewidth}{0.5pt}\end{center}

\subsection{6. Conclusion}\label{conclusion}

Dimension addition in the central projection framework is distinguished, in an interpretation-dependent manner, into two types. The conditions under which each type admits a consistent interpretation depend on the value of the parent space's curvature radius \(R\). For a flat background space (\(R = \infty\)), Type II addition is subject to the constraint of consistency from the perspective of the background space's inhabitants. For a subjective space with finite curvature (\(R < \infty\)), Type I is admissible without restriction, and Type II requires the simultaneous addition of governed dimensions. Physical interpretation lies outside the scope of this paper and is left to the reader.

\begin{center}\rule{0.5\linewidth}{0.5pt}\end{center}

\subsection{References}\label{references}

{[}1{]} Noriaki Kihara, ``Geometric Formulation via Central Projection,'' Zenodo, DOI: 10.5281/zenodo.19427780 (2025).

{[}2{]} Noriaki Kihara, ``Symmetry in the Central Projection Framework,'' Zenodo, DOI: 10.5281/zenodo.19434932 (2025).

{[}3{]} Noriaki Kihara, ``Relativity of Observation in Multiple Subjective Spaces,'' Zenodo, DOI: 10.5281/zenodo.19435162 (2025).

{[}4{]} Noriaki Kihara, ``On the Limits of Curvature Radius in Central Projection --- The Cases \(R \to \infty\) and \(R \to 0\),'' Zenodo, DOI: 10.5281/zenodo.19526549 (2026).

\begin{center}\rule{0.5\linewidth}{0.5pt}\end{center}

\end{document}
